%!TEX root = ../thesis.tex
% Worker architecture for \cref{sec:ez-loop}: the four worker types, the replay
% buffer they share, and the parameter feedback that closes the loop.
% Layout rule: the five stages form a rectangular cycle, traversed clockwise
% --- right along the top row, down the right edge, left along the bottom row,
% and up the left edge, where the parameters close it. The environment hangs
% off the cycle above the data workers, since it is not part of the algorithm.
% Every edge is horizontal or vertical: keep it that way. An earlier draft put
% the GPU workers under the environment instead of under the data workers,
% which forced the parameter arrow into a diagonal that crossed the $a_t$ edge.
% Spacing constraint: an edge label sits in the gap between two boxes, so it
% must be NARROWER than that gap or it slides under them. Columns are 5.2cm
% apart and boxes are 3cm wide, leaving a 2.2cm (62.6pt) gap; at \scriptsize
% "prioritized positions" alone is 74pt, which is why such labels are stacked
% over two lines. Widen the columns before un-stacking them. "refreshed
% targets" (61.7pt) is single-line because it spans the whole bottom row, not a
% gap, and so has 70pt of clearance on either side.
% Solid arrows carry data, dashed arrows carry parameters.
\begin{figure}[htbp]
  \centering
  \begin{tikzpicture}[
      >={Stealth[round]}, font=\small,
      worker/.style = {draw, rounded corners=2pt, fill=blue!8,
                       minimum height=9mm, minimum width=30mm, align=center},
      store/.style  = {draw, fill=black!8,
                       minimum height=9mm, minimum width=30mm, align=center},
      env/.style    = {draw, densely dashed, fill=white,
                       minimum height=9mm, minimum width=30mm, align=center},
      flow/.style   = {->, thick, black!70},
      param/.style  = {->, thick, black!70, densely dashed},
      lbl/.style    = {font=\scriptsize, inner sep=2.5pt, align=center},
    ]
    % --- the environment, outside the cycle ----------------------------------
    \node[env]    (env)  at (0,2.3)     {Environment};

    % --- top row: acting -----------------------------------------------------
    \node[worker] (data) at (0,0)       {Data workers};
    \node[store]  (buf)  at (5.2,0)     {Replay buffer};
    \node[worker] (roll) at (10.4,0)    {Rollout workers};

    % --- bottom row: learning ------------------------------------------------
    \node[worker] (gpu)  at (0,-3.1)    {GPU workers};
    \node[worker] (rean) at (10.4,-3.1) {Reanalyze workers};

    % --- environment interaction, as a parallel pair -------------------------
    \draw[flow] ([xshift=-8mm]env.south)  --
                node[lbl,left]  {$o_t,\,u_t$} ([xshift=-8mm]data.north);
    \draw[flow] ([xshift=8mm]data.north)  --
                node[lbl,right] {$a_t$}       ([xshift=8mm]env.south);

    % --- data path: clockwise around the cycle -------------------------------
    \draw[flow] (data) -- node[lbl,above] {$(o,u,\pi_t,\nu_t)$} (buf);
    \draw[flow] (buf)  -- node[lbl,above] {prioritized\\positions} (roll);
    \draw[flow] (roll) -- node[lbl,left]  {training\\batches} (rean);
    \draw[flow] (rean) -- node[lbl,above] {refreshed targets} (gpu);

    % --- parameter feedback: closes the cycle, and reaches the searchers -----
    \draw[param] (gpu) -- node[lbl,left] {$\theta$} (data);
    \draw[param,rounded corners=3pt]
      (gpu.south) -- (0,-4.1) -- node[lbl,fill=white] {$\theta$} (10.4,-4.1)
                  -- (rean.south);
  \end{tikzpicture}
  \caption[The EfficientZero worker architecture]{%
    The parallel worker architecture of EfficientZero.
    Data workers generate fresh trajectories and store them inside the replay buffer.
    Rollout and reanalyze workers sample from the replay buffer and prepare data for training.
    GPU workers calculate losses and update the model parameters.
  }
  \label{fig:ez-workers}
\end{figure}
